% Einheit 4 von 5 – Dezimalzahlen multiplizieren und dividieren (bank/bruchrechnung/e4.jsonl, zusammenbau v0.1)
%% TODO Zweigzeile Teil 1: was hier gelernt wird (Satz in Schülersprache aus dem Einheitstitel) – Einheit 4
\einheitenkopf[e4]{Einheit 4 von 5 · Dezimalzahlen multiplizieren und dividieren}
\zweigzeile{ab Kl. 5, je nach Buch bis Kl. 6 (am Gymnasium ab Kl. 6) · keine P10-Aufgabe}
\setcounter{aufgabe}{22}

%% TODO Titel: Ich-kann-Satz fehlt in der Bank (Platzhalter: Kettenname) – Nr. 23
%% TODO Anweisung: ein Satz über der ersten Teilaufgabe fehlt in der Bank – Nr. 23
\begin{aufgabe}{Wie viele Kommastellen hat das Ergebnis?}
\begin{teile}
\teil $3{,}2 \cdot 0{,}4$ – wie viele Kommastellen hat das Ergebnis? \leerfeld
\end{teile}
\end{aufgabe}

%% TODO Titel: Ich-kann-Satz fehlt in der Bank (Platzhalter: Kettenname) – Nr. 24
%% TODO Anweisung: ein Satz über der ersten Teilaufgabe fehlt in der Bank – Nr. 24
\begin{aufgabe}{Dezimal multiplizieren/dividieren}
%% TODO \rechenplatz{n}: Schrittzahl des Grundfalls fehlt in der Bank (nur bei fester Schreibform, 2.3 b)
\begin{teile}
\teil $3{,}72 \cdot 10$
\teil $45{,}7 : 10$
\teil $2{,}3 \cdot 3$
\teil $1{,}2 \cdot 0{,}7$
\teil $0{,}2 \cdot 0{,}4$
\teil $8{,}4 : 4$
\teil $5{,}6 : 0{,}7$
\teil Ein Lieferwagen braucht 8 l Diesel auf 100 km und fährt im Jahr $25\,000$ km. Ein Liter Diesel kostet $162{,}4$ ct. Berechne die Kraftstoffkosten für ein Jahr in Euro. \hfill (P10 2014 OS)
\teil Welcher Wert ist am kleinsten: $2{,}2$; $0{,}22$; $0{,}2^2$; $22\,\%$? \hfill (P10 2023 OS) \\ \leerfeld
\end{teile}
\end{aufgabe}

%% TODO Titel: Ich-kann-Satz fehlt in der Bank (Platzhalter: Kettenname) – Nr. 25
%% TODO Anweisung: ein Satz über der ersten Teilaufgabe fehlt in der Bank – Nr. 25
\begin{aufgabe}{Dezimal multiplizieren/dividieren – Fehler finden, Begründen, Anwendung}
\begin{teile}
\teil Kim rechnet: $0{,}7 \cdot 0{,}4 = 2{,}8$. Finde den Fehler.
\teil Warum rückt das Komma bei $\cdot 10$ eine Stelle nach rechts?
\teil Äpfel kosten lose 2,40 € je kg. Ein Beutel mit 1,5 kg kostet 3,29 €. Was ist günstiger?
\end{teile}
\end{aufgabe}
